% Shared preamble for the CIG-AMF formal-research papers.
% IEEEtran conference class; compiles on Overleaf without modification.
\documentclass[conference]{IEEEtran}
\IEEEoverridecommandlockouts

\usepackage{cite}
\usepackage{amsmath,amssymb,amsfonts}
\usepackage{amsthm}
\usepackage{algorithm}
\usepackage{algorithmic}
\usepackage{graphicx}
\usepackage{textcomp}
\usepackage{xcolor}
\usepackage{booktabs}
\usepackage{url}
\usepackage[hidelinks]{hyperref}

\def\BibTeX{{\rm B\kern-.05em{\sc i\kern-.025em b}\kern-.08em T\kern-.1667em\lower.7ex\hbox{E}\kern-.125emX}}

\newtheorem{theorem}{Theorem}
\newtheorem{lemma}[theorem]{Lemma}
\newtheorem{proposition}[theorem]{Proposition}
\newtheorem{corollary}[theorem]{Corollary}
\theoremstyle{definition}
\newtheorem{definition}[theorem]{Definition}
\newtheorem{remark}[theorem]{Remark}
\newtheorem{example}[theorem]{Example}

% A machine-checked statement carries its Lean declaration name.
\newcommand{\lean}[1]{\textsf{\footnotesize #1}}
\newcommand{\leanbox}[1]{%
  \par\vspace{2pt}\noindent
  {\footnotesize\textsf{Lean: #1}}\par\vspace{2pt}}

% Evidence classes, used consistently and never interchangeably.
\newcommand{\PROVED}{\textsc{proved}}
\newcommand{\CERT}{\textsc{certificate}}
\newcommand{\NUM}{\textsc{numerical}}
\graphicspath{{../_figures/}}

\usepackage{multirow}
\emergencystretch=1.5em
% DRAFT 17/10/2026. Every number traces to research/city_sim/w22_20261017 (ledger C-256..C-260); every proposition to
% work/p13_conjectures/LeanCityOwnerTestV1.lean. Bibliography details of the interference/switchback references must be verified
% against the publishers' records before submission.

\begin{document}

\title{Can a Fleet Tell That a Learned Traffic Signal Responds to It?\\
Owner-Side Switchbacks and Biased A/B Tests of Signal-Timing Services}

\author{\IEEEauthorblockN{Vu Phan An Nguyen}
\IEEEauthorblockA{\textit{Rider University} \\
Lawrenceville, NJ, USA}}

\maketitle

\begin{abstract}
A connected fleet that eco-drives on a city's signal phase and timing (SPaT) cannot see inside the city's learned signal controller,
yet whether that controller responds to the fleet's driving decides whether a timing forecast stays valid. In SUMO, with learned
signals whose inputs count only stopped vehicles, we randomise the fleet's own use of SPaT and ask two questions. First, can the fleet's
owner identify the signal's response from data it holds? A switchback over signal $\times$ 300-s windows recovers the signal's green
lengthening (MPLight $+3.85$\,s against $+4.19$\,s in the global comparison; AH $+1.69$ against $+1.51$; a threshold-free control $+0.10$
against $+0.02$), and a size-based rule decided all 50 deployed policy--fleet units correctly on fresh days when the fleet gave up SPaT in
half of the windows. With drivers matched to field trajectories part of the response depends on the network-wide share of gliding
vehicles, which a per-signal switchback cannot see and a network-wide one recovers (MPLight $+4.8$ against $+4.4$). Second, does the usual per-vehicle A/B test measure the service's harm? It does not where the signal responds: it understated
MPLight's harm by up to 2.7\,s/veh, most when few vehicles were treated, and overstated AH's at every treated share, in both driver
worlds, while it was unbiased under the threshold-free control. Running the test at two treated shares and extrapolating linearly to
all and none recovered the global effect in all five registered family--world cells, including two families and a driver world not
used to devise it. A network-wide switchback removed the bias for AH and the control but
not for MPLight, even with 1200-s windows. We also report an error in our own registered analysis: pooling green durations per signal-day is
length-biased when control windows are rare, which Chebyshev's sum inequality makes precise. The propositions are checked in Lean.
\end{abstract}

\begin{IEEEkeywords}
traffic signal control, SPaT, GLOSA, interference, switchback experiments, A/B testing, causal inference, connected vehicles, SUMO
\end{IEEEkeywords}

\section{Introduction}\label{sec:intro}
Signal phase and timing (SPaT) lets connected vehicles approach a signal at a speed that avoids a stop \cite{katsaros2011,asadi2011}.
When the signal is a learned controller, its timing is a decision taken from what it observes, and a fleet that glides changes what
it observes. Our companion study \cite{nguyen2026tsi} found that learned signals whose inputs count a vehicle only once it has stopped
(below 0.1\,m/s) serve a gliding fleet's approach later and lengthen the other greens, so the fleet's forecast defeats itself; the
inputs of sumo-rl and RESCO \cite{ault2021} are of this kind. Its practical advice, check what the signal counts, needs the signal's
internals. The fleet and the city are different owners, and the fleet's owner has only its own vehicles, the published SPaT and the
signal states its vehicles see.

This paper asks what the fleet's owner can learn from its own data. Two causal questions arise. \emph{Identification:} does the signal
respond to the fleet's use of SPaT, and by how much? \emph{Evaluation:} if the owner tests the service the usual way, by switching it on
for a random half of its vehicles, does the comparison measure what full deployment would do? The second is an interference question
\cite{savje2021,hu2022,eckles2017}: the signal reallocates green time in response to the treated vehicles, and every vehicle at the
signal receives that response. The same structure biases A/B tests in two-sided marketplaces \cite{johari2022}, in recommender
training loops \cite{si2023,brennan2023}, and in congestion-control experiments over shared links \cite{spang2021}; we found no study of it
for a responsive traffic signal, and GLOSA field evaluations compare a few vehicles with and without advice \cite{stevanovic2013}, too few
to move the signal.

\paragraph{Contributions}
\begin{enumerate}\setlength\itemsep{1pt}
\item \emph{Owner-side identification.} A switchback over signal-windows, using only owner-observable data, recovers the learned
signal's green response to within about 10\% on a 16-signal grid and detects it for the exposed families (Section~\ref{sec:d1}). A
size-based rule decides exposure without a wrong decision at a forgone SPaT share of 0.5 and 0.2; with drivers matched to field
trajectories the estimate is smaller than the global change.
\item \emph{Biased A/B tests, and a fix.} A per-vehicle A/B test misestimates the service's harm wherever the signal responds, in a
direction that depends on the signal family, follows the sign of the human drivers' change, and replicates across driver worlds; two A/B
tests at different shares, extrapolated linearly, recover the global effect; a network-wide switchback removes the bias only when the signal's
response does not build up beyond the window, which for MPLight it did even at 1200\,s (Section~\ref{sec:d2}).
\item \emph{An estimator pitfall and four propositions.} Pooling green durations per signal-day is biased downward when control windows
are rare; we found this in our own registered analysis and prove its cause. Lean also checks the A/B accounting identity, the
design-dependence of a switchback contrast under a concave exposure response, and the cost--time trade-off of the owner's test
(Section~\ref{sec:theory}).
\end{enumerate}
Every experiment was registered before it ran (protocol W22 and amendments A--D); failed predictions are listed in
Section~\ref{sec:limits}.

\section{Setting}\label{sec:setting}
We reuse the companion study's grid4 world \cite{nguyen2026tsi}: SUMO \cite{lopez2018}, 16 signals, Krauss drivers, random demand,
fresh traffic days 1400--1419, a frozen PPO residual fleet on the GLOSA rule (two of its four seeds), and the deployed SPaT service, a
learned forecast at the 0.1 quantile. Three signal families span the exposure: MPLight (RMP), whose pressure input counts halted
vehicles and is strongly exposed; the halting-count family AH, weakly exposed; and the threshold-free \texttt{wave} family (RWV), a
control. Five frozen policies per family are trained by one recipe. A replication uses the companion's W17 world, in which every car is
an IDM jointly calibrated to field trajectories, with policies, forecasts and fleets trained there. The \emph{ground truth} is the
companion's global comparison on the same policy, fleet and day: all fleet vehicles use SPaT (SV) against none (R). Mean green
lengthens by $+4.19$ (RMP), $+1.51$ (AH) and $+0.02$\,s (RWV) on grid4, and the fleet's time loss changes by $+14.05$, $+1.57$ and
$-0.96$\,s/veh.

\section{Designs, Estimands and Checks}\label{sec:design}
A wrapper randomises the fleet's \emph{use} of SPaT without editing any evaluator: the service answers every query, and a vehicle in
control ignores the answer and drives as its base driver, exactly as without SPaT. Assignments are hashes of the day and unit, so the
propensity is known.
\begin{itemize}\setlength\itemsep{1pt}
\item \textbf{MIX}$(p)$: each fleet vehicle uses SPaT with probability $p$ (the per-vehicle A/B test).
\item \textbf{SWB}$(p)$: during each 300-s window, fleet vehicles use the SPaT of signal $j$ with probability $p$, independently per
signal-window (a switchback per signal).
\item \textbf{NSW}: all fleet vehicles use SPaT in a 600-s window with probability 0.5 (network switchback); \textbf{NSWa}$(W)$: windows of
$W$\,s alternate from a random start.
\end{itemize}
The owner observes the signal states its vehicles see, hence green durations. The switchback contrast $\tau$ is the mean green
duration in treated minus control windows of a signal, over greens starting at least 60\,s into a window, averaged over signals.
A one-sided Fisher randomisation test re-draws the window labels under the sharp null that greens do not depend on the fleet's
SPaT use \cite{imbens2015}. The A/B contrast is the treated minus control fleet vehicles' mean time loss; its bias is the contrast
minus the global effect SV $-$ R.

\emph{Checks.} With every vehicle treated or none, the wrapper reproduced the SV and R rows bit for bit (time losses, fuel, counts,
teleports, the fleet's action digest and mean green) on grid4 and in the W17 world, after each edit of the wrapper. Treated shares
were $0.51$ (MIX) and $0.48$--$0.51$ (SWB). Intervals are 95\% crossed bootstraps over policies $\times$ fleets $\times$ days.

\section{Formal Results}\label{sec:theory}
All statements below are proved in Lean 4 with Mathlib (module \texttt{LeanCityOwnerTestV1}); axioms are confined to
\texttt{propext}, \texttt{Classical.choice} and \texttt{Quot.sound}, and each hypothesis set is shown satisfiable.

\begin{proposition}[Length-biased pooling]\label{prop:pool}
Let windows $i\in s$ hold $n_i>0$ greens of mean duration $m_i$. If $n$ and $m$ antivary, in particular if greens tile a window of fixed
length, $n_i=L/(m_i+c)$ with $L,c>0$, then $\sum_i n_i m_i/\sum_i n_i \le \frac{1}{|s|}\sum_i m_i$. Consequently, the contrast between
the pooled mean of the treated windows and a single control window is at most the window-weighted contrast.
\end{proposition}
\noindent The proof is Chebyshev's sum inequality. When the forgone share is $0.05$ a signal-day has on average $0.58$ control windows, so
the registered per-signal-day estimator compares one control window with about eleven pooled treated windows and is biased downward.

\begin{proposition}[A/B accounting]\label{prop:ab}
With $T_p,C_p$ the treated and control means at share $p$ and $T_1,C_0$ the all-treated and none-treated means,
$(T_p-C_p)-(T_1-C_0)=(T_p-T_1)-(C_p-C_0)$: the A/B test understates the global effect iff the treated vehicles' spillover
$T_p-T_1$ is below the control vehicles' $C_p-C_0$.
\end{proposition}

\begin{proposition}[Design-dependence]\label{prop:design}
Let a window's response be $g$ of its exposure, with own weight $a\in[0,1]$ and background $(1-a)p$, and $g$ concave on $[0,1]$. The
switchback contrast $\Delta(p)=g(a+(1-a)p)-g((1-a)p)$ is antitone in $p$, so a design with rarer control windows estimates a smaller
effect and cannot detect what a design with more control windows misses; with $g(0)=0$, $\Delta(1)\le a\,g(1)\le\Delta(0)$; a linear
response gives $\Delta\equiv\kappa a$.
\end{proposition}

\begin{proposition}[Cost and time]\label{prop:cost}
With per-day information $p(1-p)$ and an additive per-window benefit $B$, the benefit forgone until a fixed precision is
$(1-p)B\cdot c/(p(1-p))=cB/p$, decreasing in the treated share $p$, while the days needed, $c/(p(1-p))\ge 4c$, are smallest at $p=1/2$.
\end{proposition}
\noindent Proposition~\ref{prop:cost} matches the measured share of the fuel benefit the owner kept during the test, $0.47$--$0.50$,
$0.79$--$0.81$ and $0.94$--$0.96$ at $p=0.5,0.8,0.95$.

\begin{proposition}[Two-share extrapolation]\label{prop:two}
If the treated and control means are affine in the treated share, $T_p=t_0+t_1p$ and $C_p=c_0+c_1p$, then for any shares $p_1<p_2$
\[\Big(T_{p_2}+\tfrac{1-p_2}{p_2-p_1}(T_{p_2}-T_{p_1})\Big)-\Big(C_{p_1}-\tfrac{p_1}{p_2-p_1}(C_{p_2}-C_{p_1})\Big)=T_1-C_0,\]
while the naive contrast misses it by $-(1-p)t_1-p\,c_1$. If $T$ and $C$ are within $\varepsilon$ of affine at $p_1,p_2,0,1$, the
extrapolation error is at most $\varepsilon\,(4+2(1-p_2+p_1)/(p_2-p_1))$.
\end{proposition}
\noindent This is the logic of randomised-saturation designs \cite{baird2018,hudgens2008}; the propositions here and below are in module
\texttt{LeanCityOwnerTestV2}.

\begin{proposition}[What a within-run switchback cannot see]\label{prop:global}
If a window's green duration is $\mu+\tau_\ell Z+\gamma G$, with $Z$ its own assignment and $G$ the network-wide treated share, a
switchback within one run measures $\tau_\ell$, the global effect is $\tau_\ell+\gamma$, and twice the run mean over the no-SPaT baseline at
$G=1/2$ recovers $\tau_\ell+\gamma$.
\end{proposition}

\section{Results}\label{sec:results}
\subsection{Identification from the Owner's Data}\label{sec:d1}
Table~\ref{tab:d1} gives the switchback contrast against the global green change. On grid4 the owner recovered the response within
about 10\% for the exposed families, and the Fisher test detected it (MPLight $z\approx49$, AH $z\approx29$). It also detected the
threshold-free control's $+0.10$\,s ($z\approx2.9$), a response forty times smaller whose sign depended on the burn-in; the sharp-null
test detects existence, not exposure. In the W17 world the estimate was smaller than the global change for all three families under
every estimator we tried. We ruled out the metric (the same window-weighted metric on the all/none check rows gives $+6.4$ to $+6.8$\,s on
that cell), a build-up over the hour (first and last thirds $+6.30$ and $+6.33$), a lag within the window (removing the burn-in changes
nothing) and an own $\times$ neighbour interaction. Twice the run-mean green of the switchback over the baseline does recover the global
change in both worlds (W17 MPLight $+4.20$, AH $+0.61$), so in W17 about $1.0$\,s of MPLight's and $0.4$\,s of AH's response depends on the
network-wide share of gliding vehicles and is shared by treated and control windows (Proposition~\ref{prop:global}); on grid4 this part
is near zero. A network-wide switchback, registered on that basis, recovered MPLight's green change ($+4.76$ [$+2.76$, $+7.54$] against
$+4.44$; per-signal $+3.20$); for AH its interval was too wide to inform. Why the share matters with field-matched drivers is not shown;
one candidate is platoon dispersion from upstream gliding that persists across blocks.

\begin{table}[t]\centering\caption{Owner-side switchback (SWB, $p=0.5$): green contrast $\tau$ (s) against the global change SV $-$ R}\label{tab:d1}
\footnotesize\resizebox{\columnwidth}{!}{\begin{tabular}{@{}llrrr@{}}\toprule
World & Family & $\tau$ (registered) & $\tau$ (window-wtd.) & SV $-$ R \\\midrule
grid4 & MPLight & $+3.85$ [$+2.46$, $+5.59$] & $+3.97$ & $+4.19$ \\
grid4 & AH & $+1.69$ [$+1.16$, $+2.40$] & $+1.79$ & $+1.51$ \\
grid4 & \texttt{wave} & $+0.10$ [$-0.12$, $+0.32$] & $+0.10$ & $+0.02$ \\
W17 & MPLight & $+3.20$ [$+1.97$, $+4.74$] & $+3.24$ & $+4.44$ \\
W17 & AH & $+0.16$ [$-0.33$, $+0.62$] & $+0.18$ & $+0.66$ \\
W17 & \texttt{wave} & $-0.28$ [$-0.53$, $-0.06$] & $-0.28$ & $-0.02$ \\\bottomrule
\end{tabular}}\end{table}

\emph{Deciding by size.} On fresh days 1410--1419, ten policy--fleet units per family were labelled \emph{exposed} once a Student-$t$
interval of the daily contrasts lay above $\delta=0.5$\,s, and \emph{not exposed} once it lay below (Table~\ref{tab:dec}). With the
registered estimator, giving up SPaT in half the windows decided all 50 units correctly in a median of 3 days (at most 7); at a
forgone share of $0.2$ there was again no wrong decision; at $0.05$ most units stayed undecided and wrong decisions appeared. The
apparent shrinkage of $\tau$ with the treated share (MPLight $3.63\to2.63$) is mostly the bias of Proposition~\ref{prop:pool}: with the
window-weighted estimator (post hoc), $\tau$ is nearly flat ($3.83$, $3.92$, $3.48$) and one wrong decision remains at $0.05$.

\begin{table}[t]\centering\caption{Size-based decisions (50 units; days 1410--1419): wrong / undecided / range of the families' median days}\label{tab:dec}
\footnotesize\resizebox{\columnwidth}{!}{\begin{tabular}{@{}lccc@{}}\toprule
Forgone share & 0.5 & 0.2 & 0.05 \\\midrule
Registered estimator & 0 / 0 / 3--4 & 0 / 3 / 3--6.5 & 1 / 26 / 2--6.5 \\
Window-weighted (post hoc) & 0 / 0 / 2.5--5 & 0 / 1 / 3--5 & 1 / 23 / 4--9 \\
Fuel benefit kept & 0.48 & 0.80 & 0.95 \\\bottomrule
\end{tabular}}\end{table}

\subsection{Per-Vehicle A/B Tests}\label{sec:d2}
Table~\ref{tab:ab} gives the bias of the per-vehicle contrast. Under MPLight the test understated the harm at every treated share,
most when few vehicles were treated ($-2.74$\,s/veh at $p=0.1$, about a fifth of the harm): a few gliding vehicles do not move the
signal, so a small test misses the feedback. Under AH it overstated the harm at every share ($+1.25$ to $+1.59$): control vehicles gained
from the treated ones' gliding ($C_p-C_0$ from $-0.13$ to $-1.45$), which Proposition~\ref{prop:ab} turns into an overstatement. Under the
threshold-free control no bias was shown. In the W17 world the signs replicated (MPLight $-1.28$, AH $+0.67$, \texttt{wave} $-0.05$).

\begin{table}[t]\centering\caption{Per-vehicle A/B bias (s/veh) against SV $-$ R; grid4, days 1400--1409}\label{tab:ab}
\footnotesize\begin{tabular}{@{}lrrrrr@{}}\toprule
Treated share & 0.1 & 0.3 & 0.5 & 0.7 & 0.9 \\\midrule
MPLight & $-2.74^*$ & $-2.34^*$ & $-1.97^*$ & $-1.28^*$ & $-0.65$ \\
AH & $+1.46^*$ & $+1.29^*$ & $+1.25^*$ & $+1.29^*$ & $+1.59^*$ \\
\texttt{wave} & $+0.24$ & $+0.12$ & $+0.07$ & $+0.26$ & $+0.78$ \\\bottomrule
\multicolumn{6}{@{}l}{\scriptsize $^*$ interval excludes 0}
\end{tabular}\end{table}

\emph{The direction of the bias.} Proposition~\ref{prop:ab} splits the bias into the treated and control vehicles' spillovers. Under
MPLight the treated vehicles fare better in a mixed world than when all glide ($T_p-T_1$ from $-2.65$ at $p=0.1$ to $-0.20$ at $0.9$): the
feedback needs many gliding vehicles, so a small test misses it. Under AH they fare worse ($+1.33$ to $+0.14$) while control vehicles gain:
a signal that serves halted vehicles first lets the counted vehicles go ahead of the uncounted ones. The control vehicles' spillover had the
sign of the human drivers' change under the service in every registered cell (grid4 MPLight, AH, SRL, RDQ; W17 MPLight, AH), so a city that
sees its drivers gain when the fleet uses SPaT can expect the fleet's A/B test to overstate the harm. As registered, the naive test
overstated it for RDQ ($+1.24$ [$+0.63$, $+1.94$]); for SRL the sign was right but the interval reached zero ($+0.49$ [$-0.04$, $+1.03$]).

\emph{Two shares.} Because $T_p$ and $C_p$ were close to linear in $p$, we registered Proposition~\ref{prop:two}'s extrapolation from
shares $0.3$ and $0.7$ after seeing it work post hoc on Table~\ref{tab:ab}'s runs (MPLight $14.51$ against $14.05$; AH $1.38$ against $1.57$).
On two families and a driver world not used to devise it, its bias interval contained zero in all five cells (Table~\ref{tab:two}). Our
two shares were run on the same simulated days, so its precision here is optimistic.

\begin{table}[t]\centering\caption{Two-share extrapolation (0.3, 0.7) against the global effect (s/veh); registered cells}\label{tab:two}
\footnotesize\resizebox{\columnwidth}{!}{\begin{tabular}{@{}llrrr@{}}\toprule
World & Family & GTE & two-share bias & naive A/B bias \\\midrule
grid4 & SRL & $+4.64$ & $+0.00$ [$-0.81$, $+0.80$] & $+0.49$ [$-0.04$, $+1.03$] \\
grid4 & RDQ & $+2.38$ & $+0.09$ [$-0.42$, $+0.65$] & $+1.24$ [$+0.63$, $+1.94$] \\
W17 & MPLight & $+14.05$ & $-0.04$ [$-0.68$, $+0.56$] & $-1.28$ [$-2.34$, $-0.14$] \\
W17 & AH & $+1.81$ & $+0.02$ [$-0.39$, $+0.43$] & $+0.67$ [$+0.12$, $+1.22$] \\
W17 & \texttt{wave} & $+0.92$ & $+0.19$ [$-0.22$, $+0.57$] & $-0.05$ [$-0.58$, $+0.46$] \\\bottomrule
\end{tabular}}\end{table}

\emph{Whole days.} Randomising whole days between the service and no service is unbiased by construction; with the day-to-day spread of our
simulated days (0.5--0.65\,s/veh) a unit needs about 2--4 days for AH and 8--10 for the control at 80\% power, but real days vary more.

\emph{Network switchbacks.} Switching the service for the whole fleet in 600-s windows recovered the global effect for AH (bias $-0.34$
[$-1.12$, $+0.46$]) and the control ($+0.14$), but not for MPLight ($-1.85$ [$-3.14$, $-0.63$]). Longer alternating windows did not help: MPLight's bias was $-2.71$
[$-4.36$, $-1.34$] with 900-s and $-1.96$ [$-3.51$, $-0.68$] with 1200-s windows, and AH's became $-0.71$ [$-1.38$, $-0.09$] at 1200\,s,
while the control stayed unbiased. The start of the alternation was balanced over days (5 / 5), so a time trend within the hour does not
explain it; we read it as a response that builds up over the whole hour, queues growing under lengthened greens, which windows of 10--20
minutes never reach. For such a controller the owner's evaluation needs randomisation over whole days or areas, which we have not tested.

\section{Related Work}\label{sec:related}
Interference makes a unit's outcome depend on others' treatments, so the estimand depends on the design
\cite{savje2021,hu2022}; cluster and graph randomisation reduce the bias \cite{eckles2017}. Switchback designs randomise over time and
must handle carry-over \cite{bojinov2023}. Two-stage and randomised-saturation designs vary the treated share across clusters to separate direct from spillover
effects \cite{hudgens2008,baird2018}; our two-share extrapolation applies that idea over days. Marketplace experiments show that the bias of customer- and listing-side randomisation
depends on market balance \cite{johari2022}, and shared training pipelines bias A/B tests of recommenders \cite{si2023,brennan2023}. The
closest structure is A/B testing of congestion control over a shared bottleneck, where tests misjudged the direction of some metrics
\cite{spang2021}. Here the shared resource is green time and the interference runs through a learned controller's sensing. Performative
prediction studies forecasts that change their target \cite{perdomo2020,narang2023,mandal2023}; the companion study found such a loop
for SPaT \cite{nguyen2026tsi}. GLOSA evaluations under actuated signals report unreliable timing \cite{stevanovic2013} but use
within-vehicle comparisons.

\section{Limitations and Failed Predictions}\label{sec:limits}
\begin{itemize}\setlength\itemsep{1pt}
\item \emph{Simulation only}, one 16-signal grid and one field-matched world; fleets, signals and forecasts are frozen policies from one
recipe. Only the fleet's \emph{use} of SPaT is randomised; real deployments add detector errors and partial penetration.
\item \emph{The W17 network-share component} is measured, not explained; MPLight's network-switchback bias on grid4 is not explained either.
The two-share extrapolation assumes near-linearity in the share and was tested with the two shares on the same simulated days.
\item \emph{The threshold $\delta=0.5$\,s} was chosen after the first panel had been seen; it was evaluated only on fresh days and in the
W17 world. The confidence sequence uses a plug-in variance and is approximate.
\item \emph{Registered predictions that failed:} the threshold-free control was detected; ``every exposed unit within 3 days'' (the
maximum was 7); ``no network-switchback bias for any family'' (MPLight); ``MPLight's bias shrinks with longer windows'' and ``AH unbiased at 900 and 1200\,s'' (AH at 1200\,s); ``the naive test overstates SRL's harm'' (right sign, interval reaching zero). The claim that cheap tests are biased through carry-over, which
we first recorded, was withdrawn: it was the estimator bias of Proposition~\ref{prop:pool}.
\item \emph{The propositions} are deterministic statements about the models stated; that SUMO windows tile exactly and that the real
exposure response is concave are not established, and the corrected estimates are nearly flat in the treated share.
\end{itemize}

\section{Conclusion}\label{sec:concl}
A fleet's owner does not need the city's controller to learn whether it responds: randomising its own vehicles' use of SPaT over
signal-windows recovers the response and supports a decision at a cost it can choose. The same responsiveness makes the obvious
evaluation, a per-vehicle A/B test, unreliable in a direction that depends on the controller. Before trusting a timing service's
measured effect, randomise at the level at which the controller responds, or test at two shares and extrapolate, and pool by window
rather than by green.

\begin{thebibliography}{00}\footnotesize
\bibitem{asadi2011} B.~Asadi, A.~Vahidi. Predictive cruise control: Utilizing upcoming traffic signal information for improving fuel economy and reducing trip time. \emph{IEEE Trans. Control Syst. Technol.} 19(3):707--714, 2011.
\bibitem{ault2021} J.~Ault, G.~Sharon. Reinforcement learning benchmarks for traffic signal control. \emph{NeurIPS Datasets and Benchmarks Track}, 2021.
\bibitem{baird2018} S.~Baird, J.~A.~Bohren, C.~McIntosh, B.~\"Ozler. Optimal design of experiments in the presence of interference. \emph{Rev. Econ. Stat.} 100(5):844--860, 2018.
\bibitem{bojinov2023} I.~Bojinov, D.~Simchi-Levi, J.~Zhao. Design and analysis of switchback experiments. \emph{Management Science} 69(7):3759--3777, 2023.
\bibitem{brennan2023} J.~Brennan, V.~Mirrokni, J.~Pouget-Abadie. A study of symbiosis bias in A/B tests of recommendation algorithms. arXiv:2309.07107, 2023.
\bibitem{eckles2017} D.~Eckles, B.~Karrer, J.~Ugander. Design and analysis of experiments in networks: Reducing bias from interference. \emph{J. Causal Inference} 5(1), 2017.
\bibitem{hudgens2008} M.~G.~Hudgens, M.~E.~Halloran. Toward causal inference with interference. \emph{J. Am. Stat. Assoc.} 103(482):832--842, 2008.
\bibitem{hu2022} Y.~Hu, S.~Li, S.~Wager. Average direct and indirect causal effects under interference. \emph{Biometrika} 109(4):1165--1172, 2022.
\bibitem{imbens2015} G.~W.~Imbens, D.~B.~Rubin. \emph{Causal Inference for Statistics, Social, and Biomedical Sciences}. Cambridge University Press, 2015.
\bibitem{johari2022} R.~Johari, H.~Li, I.~Liskovich, G.~Y.~Weintraub. Experimental design in two-sided platforms: An analysis of bias. \emph{Management Science} 68(10):7069--7089, 2022.
\bibitem{katsaros2011} K.~Katsaros, R.~Kernchen, M.~Dianati, D.~Rieck. Performance study of a Green Light Optimized Speed Advisory (GLOSA) application using an integrated cooperative ITS simulation platform. \emph{Proc. IWCMC}, 2011.
\bibitem{lopez2018} P.~A.~Lopez et al. Microscopic traffic simulation using SUMO. \emph{Proc. IEEE ITSC}, 2018.
\bibitem{mandal2023} D.~Mandal, S.~Triantafyllou, G.~Radanovic. Performative reinforcement learning. \emph{Proc. ICML}, 2023.
\bibitem{narang2023} A.~Narang, E.~Faulkner, D.~Drusvyatskiy, M.~Fazel, L.~J.~Ratliff. Multiplayer performative prediction: Learning in decision-dependent games. \emph{J. Mach. Learn. Res.} 24(202), 2023.
\bibitem{nguyen2026tsi} V.~P.~A.~Nguyen. When gliding cars leave the queue: Learned traffic signals that count stopped vehicles defeat eco-driving fleets, and repairs without retraining. Manuscript, 2026.
\bibitem{perdomo2020} J.~Perdomo, T.~Zrnic, C.~Mendler-D\"unner, M.~Hardt. Performative prediction. \emph{Proc. ICML}, PMLR 119:7599--7609, 2020.
\bibitem{savje2021} F.~S\"avje, P.~M.~Aronow, M.~G.~Hudgens. Average treatment effects in the presence of unknown interference. \emph{Ann. Statist.} 49(2):673--701, 2021.
\bibitem{si2023} N.~Si. Tackling interference induced by data training loops in A/B tests: A weighted training approach. arXiv:2310.17496, 2023.
\bibitem{spang2021} B.~Spang, V.~Hannan, S.~Kunamalla, T.-Y.~Huang, N.~McKeown, R.~Johari. Unbiased experiments in congested networks. \emph{Proc. ACM IMC}, 2021.
\bibitem{stevanovic2013} A.~Stevanovic, J.~Stevanovic, C.~Kergaye. Green light optimized speed advisory systems: Impact of signal phasing information accuracy. \emph{Transp. Res. Rec.} 2390:53--59, 2013.
\end{thebibliography}

\end{document}
